% Määrittele avainsanat main.tex:ssä ja poista seuraava rivi
\avainsanat{diplomi-insinöörin tutkinto\spc diplomityön ohjeet\spc diplomityön rakenne\footnote{Käytä avainsanoina eri sanoja kuin lopputyön otsikossa. Poista tämä alaviite lopputyöstäsi.}}
\begin{thesisabstract}{finnish}
\tiivistelmametadata

Tässä työssä kehitetään suoritettava kehys indonesialaisten
vastuullisuusraporttien analysointiin tietuepohjaisena evidenssinä sen sijaan,
että raportteja käsiteltäisiin yhtenäisinä kokonaisuuksina. Työ ratkaisee
teknisen ongelman, jossa heterogeeniset ja kaksikieliset PDF-muotoiset raportit
muunnetaan rakenteisiksi ja auditoitaviksi Environmental, Social and
Governance -sisältöä kuvaaviksi tietueiksi. Toteutettu työnkulku yhdistää
OCR-pohjaisen dokumenttien laajennuksen, sivutasoisen alkuperän
jäljitettävyyden, suurten kielimallien avulla tehdyn poiminnan,
tietuetasoisen skeeman, ontologiaan kohdistamisen sekä vertailun
ClimateBERT-tyyppisiin ilmastoluokituksiin.

Kehystä arvioitiin työn aktiivisella osajoukolla, joka sisälsi 23 käsiteltyä
raporttia ja noin 5 512 sivua. Putki tuotti 332 sävyluokiteltua ESG-tietuetta
ja 2 074 jatkoanalyysin riviä, ja kaikki 52 seurattua aspektia pystyttiin
liittämään ontologiapolkuihin. Tulokset osoittavat, että ilmaisun sävy on
mitattava ja hyödyllinen: sitoumussävy oli yleisin, ympäristö oli yleisin
ESG-ulottuvuus, ja sävyluokat saavuttivat 83,7 prosentin yhtäpitävyyden
ClimateBERT-tyyppisten sitoumusluokkien kanssa Cohenin kappan ollessa 0,645.

Samalla kokeet osoittavat, että poiminnan laatu riippuu vahvasti kehotteiden
suunnittelusta ja mallin skeemanoudatuksesta, ja suurimmat heikkoudet liittyvät
puuttuviin sävyluokkiin ja skeemavirheisiin. Työn johtopäätös on, että
sävy­tietoinen ja jäljitettävyyden säilyttävä ESG-analyysi on
toteuttamiskelpoinen ja informatiivinen, mutta vahvempi asiantuntija-annotointi,
laajempi kontrolloitu validointi ja muodolliset ablaatiokokeet ovat edelleen
tarpeen ennen benchmark-tasoisten suorituskykyväitteiden esittämistä.
\end{thesisabstract}
